\documentclass[./book.tex]{subfiles}

\begin{document}

    \section{计算机系统概述}

    \begin{table}[htbp]
        \centering
        \renewcommand{\arraystretch}{1.5}   % 行高
        \setlength{\tabcolsep}{6pt}         % 列间距
        \small
        \begin{tabular}{|
                >{\centering\arraybackslash}m{1.8cm}|
                >{\centering\arraybackslash}m{3.5cm}|
                >{\centering\arraybackslash}m{10cm}|}
            \hline
            \textbf{名称} & \textbf{英文}         & \textbf{含义}                              \\
            \hline
            存储元         & Memory Cell         & 一个二进制位（1 bit）                            \\
            \hline
            存储单元        & Storage Unit        & 一个地址对应的数据                                \\
            \hline
            存储字         & Memory Word         & 一个存储单元存放的数据（CPU一次存取的数据单位、存储单元存储的一串二进制代码） \\
            \hline
            存储字长        & Memory Word Length  & 一个存储单元包含的二进制位数（由编址方式决定）                  \\
            \hline
            机器字长        & Machine Word Length & CPU 一次能够自然处理的二进制数据位数（由 CPU 决定）           \\
            \hline
        \end{tabular}
    \end{table}

    \begin{enumerate}
        \item 按地址存取方式：（主存储器的工作方式是）按存储单元的地址进行存取
        \item MAR（存储器地址寄存器，Memory Address Register）：存放访存地址，位数反映最多可寻址的存储单元的个数，位数与地址线位数相同
        \begin{itemize}
            \item 地址个数 = 终止地址 - 起始地址 + 1
            \item 地址位数 = $\log_2 \textbf{存储单元个数}$
            \begin{circlenum}
                \item 字节编址（存储字长1B = 8bit）：存储单元个数 = 存储器容量/1B。\emph{e.g.} 主存容量$32MB$，\textbf{字节}编址 $\Rightarrow$ 地址位数$= 2^5 * 2^{20}$
                \item 32位字编址（存储字长32bit）：存储单元个数 = 存储器容量/32bit。\emph{e.g.} 主存容量$32MB$，\textbf{字}编址 $\Rightarrow$ 地址位数$= 2^{20}$
            \end{circlenum}
        \end{itemize}
        \item MDR（存储器数据寄存器，Memory Data Register）：用于暂存要从存储器中读或写的信息，位数与数据总线宽度（位数）相同
        \begin{analysisbox}[Example]
            某 64 位计算机，主存容量 4MB，字节编址 \\
            $\Rightarrow$ \begin{cases*}
                              \text{机器字长} = 64 bit \\
                              \text{存储字长} = 1B = 8 bit \\
                              MDR\textbf{位数} = 64 bit \\
                              MAR\textbf{位数} = \log_2 \dfrac{4MB}{1B} = 2^{22} \\
            \end{cases*}
        \end{analysisbox}
        \item CPU 一次从主存读/写多少数据，取决于：
        \begin{itemize}
            \item 数据总线宽度
            \item 存储器接口设计
            \item 一次访存操作的宽度
            \item 数据总线：64bit，那么一次访存可以读： $64bit=8Byte$。也就是一次可以读 8 个按字节编址的存储单元
        \end{itemize}
    \end{enumerate}

    \subsection{计算机的性能指标}

    \begin{enumerate}
        \item CPU时钟周期：CPU工作的最小时间单位
        \item 主频（CPU时钟频率）：每秒有多少个时钟周期，单位$Hz$
        \[
            \text{主频} = \dfrac{1}{\text{CPU时钟周期}}
        \]
        \item CPI（CyclePerInstruction）：执行一条指令所需的时钟周期数
        \item IPS（Instruction Per Second）：每秒执行多少条指令
        \[
            IPS = \dfrac{\text{主频}}{\text{平均}CPI}
        \]
        \item CPU执行时间：运行一个程序所花费的时间
        \[
            CPU\text{执行时间} = \dfrac{CPU\text{时钟周期数}}{\text{主频}} = \dfrac{\text{指令条数}\times CPI}{\text{主频}}
        \]
        \item MIPS（Million Instruction Per Second）：每秒执行多少百万条指令
        \[
            MIPS = \dfrac{\text{指令条数}}{\text{执行时间}\times 10^6} = \dfrac{\text{主频}}{CPI * 10^6}
        \]
        \item 机器的速度与基准程序在该机器上的运行时间呈相反关系
    \end{enumerate}

\end{document}
